\documentclass{article}
\usepackage{enumitem}
\usepackage{amsmath}
\begin{document}

\title{Chapter 23: Reproduction in Organisms}
\date{}
\maketitle

\begin{enumerate}[label=Q\arabic*.]

\item Which of the following is NOT an advantage of sexual reproduction over asexual reproduction?
\\(A) Genetic variation in offspring
\\(B) Faster rate of reproduction
\\(C) Adaptation to changing environments
\\(D) New gene combinations through recombination

\item Vegetative propagation in plants is a form of:
\\(A) Sexual reproduction
\\(B) Asexual reproduction producing genetically identical offspring (clones)
\\(C) Amphimixis
\\(D) Parthenogenesis

\item In binary fission of \textit{Amoeba}, the process produces:
\\(A) Two unequal daughter cells
\\(B) Two genetically identical daughter cells
\\(C) Four haploid cells
\\(D) Spores that germinate

\item Budding as seen in \textit{Hydra} involves:
\\(A) Meiotic division of germ cells
\\(B) Mitotic outgrowth from the parent body that eventually detaches as a new individual
\\(C) Fragmentation followed by regeneration
\\(D) Sporulation

\item Assertion: Offspring produced by asexual reproduction are genetically identical to the parent.\\
Reason: Asexual reproduction involves only mitosis without fusion of gametes.
\\(A) Both true, Reason is correct explanation
\\(B) Both true, Reason is not correct explanation
\\(C) Assertion true, Reason false
\\(D) Assertion false, Reason true

\item Parthenogenesis is defined as:
\\(A) Development of embryo from an unfertilized egg
\\(B) Fusion of two identical gametes
\\(C) Reproduction without a partner by vegetative means
\\(D) Formation of seeds without pollination

\item Which of the following organisms reproduces by sporulation?
\\(A) \textit{Hydra}
\\(B) \textit{Rhizopus} (bread mould)
\\(C) \textit{Planaria}
\\(D) \textit{Amoeba}

\item The term ``juvenile phase'' (vegetative phase in plants) refers to:
\\(A) The phase after an organism reaches reproductive maturity
\\(B) The phase from birth/germination to the onset of reproductive capacity
\\(C) The dormancy period in seeds
\\(D) The larval stage in insects

\item Senescence in plants is characterized by:
\\(A) Increased growth rate
\\(B) Progressive deterioration of cellular functions leading to death, associated with ethylene and ABA
\\(C) Increased photosynthetic rate
\\(D) Accumulation of cytokinins

\item Which of the following is an example of external fertilization?
\\(A) Humans
\\(B) Frogs and fish releasing eggs and sperm into water
\\(C) Birds
\\(D) Reptiles

\item Oviparous animals differ from viviparous animals in that:
\\(A) Oviparous animals always have internal fertilization
\\(B) Oviparous animals lay eggs; embryo develops outside the mother's body
\\(C) Oviparous animals provide more parental care
\\(D) Oviparous animals produce fewer offspring

\item Regeneration as seen in \textit{Planaria} involves:
\\(A) Meiosis to produce new individuals
\\(B) Ability of a fragment to grow into a complete organism via mitosis and differentiation
\\(C) Fusion of gametes
\\(D) Asexual spore formation

\item In \textit{Rhizopus}, the sporangiospores are:
\\(A) Sexual spores formed by meiosis
\\(B) Asexual spores formed by mitosis within sporangia
\\(C) Formed by budding
\\(D) Produced by fragmentation

\item Match the types of reproduction with examples:\\
1. Binary fission $\rightarrow$ (a) \textit{Amoeba}, bacteria\\
2. Budding $\rightarrow$ (b) \textit{Hydra}, yeast\\
3. Fragmentation $\rightarrow$ (c) \textit{Spirogyra}\\
4. Sporulation $\rightarrow$ (d) \textit{Rhizopus}
\\(A) 1-a, 2-b, 3-c, 4-d
\\(B) 1-b, 2-a, 3-d, 4-c
\\(C) 1-c, 2-d, 3-a, 4-b
\\(D) 1-d, 2-c, 3-b, 4-a

\item The biological term for the lifespan between birth and natural death is:
\\(A) Senescence
\\(B) Longevity
\\(C) Reproductive phase
\\(D) Juvenile phase

\end{enumerate}
\end{document}
